\chapter{Conclusion and Work Plan}\label{ch:concl}

This report extends the previous grouped-MAPPO study from communication-module
comparison to the joint design of representation, grouping, and routing. The
main architectural contribution is TarMAC Hybrid. Its ReLU fusion compresses a
256-dimensional local state to 64 dimensions, combines that representation with
a 64-dimensional attention context, and adds a bounded residual update. In the
matched initial benchmark, it lowers CV-RMSE from 49.65\% to 47.01\% and comfort
exceedance from 22.10\% to 17.99\% relative to the original TarMAC adaptation.

\begin{zhtranslation}
本报告把上一阶段 grouped MAPPO 的通信模块比较，扩展为表示、分组和路由的联合设计。
主要结构贡献是 TarMAC Hybrid：ReLU fusion 把 256 维本地状态压缩到 64 维，与 64 维
注意力上下文结合，再施加有界残差更新。在匹配的初始基准中，它相对原始 TarMAC 把
CV-RMSE 从 49.65\% 降到 47.01\%，把舒适超限从 22.10\% 降到 17.99\%。
\end{zhtranslation}

The grouping study shows that neither the algorithm nor the features can be
chosen independently of policy learning. Agglomerative clustering is the most
useful deterministic baseline for the subsequent feature study, while balanced
spectral produces the best individual compact-group run but remains seed-
sensitive. The feature-count ablation shows that 4F and 5F can be worse than a
compact 3F description. Among the three-feature candidates, 3fA is selected as
the overall compromise because it provides extremely low and stable portfolio
bias, preserves distinct physical roles, and transfers naturally from Vermont
heating to Texas cooling. It is not universally best: 3fI is stronger on some
mean shape and comfort metrics, which should remain visible in the final thesis.

\begin{zhtranslation}
分组研究表明，聚类算法和特征都不能脱离策略学习单独选择。凝聚聚类是后续特征实验最
有用的确定性基线；balanced spectral 产生最佳单次 compact-group 结果，但仍对 seed
敏感。特征数量消融表明 4F 和 5F 可能差于紧凑 3F。在三特征候选中，3fA 因组合 bias
极低且稳定、物理角色清晰、并能从 Vermont 供暖自然迁移到 Texas 制冷，被选为综合
折中。但它并非普遍第一：3fI 在部分平均负荷形状和舒适指标上更强，最终论文应保留
这一事实。
\end{zhtranslation}

The soft-router study gives a similarly qualified positive result. The stable
full-expert/head-adaptation configuration obtains CV-RMSE
$45.12\pm0.35$\% over three router seeds that share one seed-42 expert
checkpoint. It is therefore stable conditional on a strong anchor, rather than
a complete end-to-end stability result. The router does not learn useful
organization from scratch: legacy end-to-end and shared three-stage training
are weaker, while strong results appear only when a good fixed-group checkpoint
is preserved. Soft routing is a valuable adaptive enhancement, but not yet a
replacement for grouping.

\begin{zhtranslation}
Soft router 研究也得到带条件的正面结论。稳定 full-expert/head-adaptation 配置在三个
router seeds 上取得 $45.12\pm0.35$\% 的 CV-RMSE，但三次都共享同一个 seed-42 expert
checkpoint，因此这是强锚点条件下的稳定性，不是完整端到端稳定性。Router 不能从零学出
有效组织：legacy end-to-end 和 shared 三阶段较弱，只有保留优质固定分组 checkpoint
时才出现强结果。因此，soft routing 是有价值的自适应增强，但尚不能替代分组。
\end{zhtranslation}

\section{Next Steps}

The remaining work should focus on validation rather than adding many new
architectures:
\begin{itemize}
  \item run a paired seed set for fixed 3fA, router-only, and stable-head routing
  so that improvements can be tested seed by seed;
  \item investigate the anomalous stable-head reward in seed 1 and verify that
  reward aggregation, comfort penalty, and exported primary metrics use the
  same evaluation interval;
  \item repeat the stable router in Texas with cooling-aware dynamic inputs and
  compare against the fixed TX 3fA controller;
  \item report cluster assignments and gate-usage heat maps to show which
  buildings actually switch experts and under which load/SOC conditions;
  \item test whether a weak initial partition can be repaired by the router, or
  whether performance is monotonic in fixed-expert quality;
  \item retain 3fI as a competing feature set and evaluate whether its lower
  comfort error persists in paired multi-seed and TX experiments.
\end{itemize}

\begin{zhtranslation}
后续工作应优先加强验证，而不是继续增加大量新结构：对固定 3fA、router-only 和
stable-head 使用完全配对的 seeds；调查 stable-head seed 1 的异常 reward，并确认 reward
聚合、舒适度惩罚和主指标使用同一测试区间；在 Texas 中用 cooling-aware 动态输入复现
稳定 router，并与固定 TX 3fA 比较；报告 cluster assignment 和 gate usage heat map，
说明哪些建筑在何种负荷/SOC 条件下切换 experts；测试 router 能否修复较弱初始分组，
或性能是否始终受固定 expert 质量限制；保留 3fI 作为竞争特征集，检验其舒适度优势在
配对多种子和 TX 中是否继续存在。
\end{zhtranslation}
